\documentclass[11pt]{article}
\usepackage{amsmath, amssymb, amsthm, mathtools, hyperref}

\newtheorem{theorem}{Theorem}
\newtheorem{definition}{Definition}
\newtheorem{lemma}{Lemma}
\newtheorem{proposition}{Proposition}

\title{Logic and Proof}
\author{math-foundations / 01-logic-and-proof}
\date{}

\begin{document}
\maketitle

\iffalse
This lesson covers propositional logic, predicate logic, truth tables, and the
four classical proof strategies, with a full inductive proof of the closed form
for the first $n$ positive integers.
\fi

\section*{Introduction}

A \emph{proposition} is a declarative sentence with a definite truth value in
$\{\text{True}, \text{False}\}$. A \emph{proof} is a finite chain of
propositions, each justified by an axiom or a prior theorem under valid rules
of inference. Every theoretical claim in mathematics --- and every formal
guarantee in a machine-learning paper --- ultimately reduces to such a chain.
The goal of this lesson is to make the vocabulary precise and to drill the four
basic proof strategies.

\section*{Definitions}

\begin{definition}[Proposition]
A \emph{proposition} is a statement that is either true or false (but not
both). Variables $p, q, r$ range over propositions.
\end{definition}

\begin{definition}[Logical connectives]
For propositions $p, q$ we define:
\begin{align*}
\neg p &: \text{``not } p\text{''}, \quad \text{true iff } p \text{ is false}; \\
p \land q &: \text{``}p \text{ and } q\text{''}, \quad \text{true iff both are true}; \\
p \lor q &: \text{``}p \text{ or } q\text{''}, \quad \text{true iff at least one is true}; \\
p \Rightarrow q &: \text{``if } p \text{ then } q\text{''}, \quad \text{false iff } p \text{ true and } q \text{ false}; \\
p \Leftrightarrow q &: \text{``}p \text{ iff } q\text{''}, \quad \text{true iff } p \text{ and } q \text{ have the same value}.
\end{align*}
\end{definition}

\begin{definition}[Tautology]
A compound proposition is a \emph{tautology} if it is true under every
assignment of truth values to its atomic variables.
\end{definition}

\begin{definition}[Predicate, quantifiers]
A \emph{predicate} $P(x)$ on a domain $D$ is a statement whose truth value
depends on $x \in D$. The \emph{universal quantifier} $\forall x \in D,\, P(x)$
asserts $P$ holds for every $x$; the \emph{existential quantifier}
$\exists x \in D,\, P(x)$ asserts $P$ holds for at least one $x$.
\end{definition}

\section*{Properties and Theorems}

\begin{proposition}[Contrapositive equivalence]
For all propositions $p, q$, $(p \Rightarrow q) \Leftrightarrow (\neg q \Rightarrow \neg p)$.
\end{proposition}

\begin{proposition}[De Morgan]
$\neg(p \land q) \equiv (\neg p) \lor (\neg q)$ and
$\neg(p \lor q) \equiv (\neg p) \land (\neg q)$. For quantifiers,
$\neg \forall x\, P(x) \equiv \exists x\, \neg P(x)$ and
$\neg \exists x\, P(x) \equiv \forall x\, \neg P(x)$.
\end{proposition}

\begin{theorem}[Principle of Mathematical Induction]
\label{thm:induction}
Let $P(n)$ be a predicate on $\mathbb{N} = \{1, 2, 3, \dots\}$. Suppose
\begin{enumerate}
  \item (Base case) $P(1)$ is true, and
  \item (Inductive step) for every $k \in \mathbb{N}$, $P(k) \Rightarrow P(k+1)$.
\end{enumerate}
Then $P(n)$ is true for every $n \in \mathbb{N}$.
\end{theorem}

\section*{Worked Examples}

\subsection*{Direct proof}

\begin{proposition}
If $n$ is an even integer, then $n^2$ is even.
\end{proposition}
\begin{proof}
By hypothesis $n = 2k$ for some $k \in \mathbb{Z}$. Then
$n^2 = (2k)^2 = 4k^2 = 2(2k^2)$, which is twice an integer, hence even.
\end{proof}

\subsection*{Contrapositive}

\begin{proposition}
If $n^2$ is an even integer, then $n$ is even.
\end{proposition}
\begin{proof}
We prove the contrapositive: if $n$ is odd, then $n^2$ is odd. Write
$n = 2k+1$. Then $n^2 = 4k^2 + 4k + 1 = 2(2k^2 + 2k) + 1$, which is odd.
\end{proof}

\subsection*{Contradiction}

\begin{proposition}
$\sqrt{2}$ is irrational.
\end{proposition}
\begin{proof}
Suppose for contradiction $\sqrt{2} = a/b$ with $a, b \in \mathbb{Z}$, $b \neq 0$,
$\gcd(a,b) = 1$. Then $a^2 = 2b^2$, so $a^2$ is even, so $a$ is even
(contrapositive lemma above). Write $a = 2c$. Then $4c^2 = 2b^2$, so
$b^2 = 2c^2$, so $b$ is even. But then $\gcd(a,b) \geq 2$, contradicting
$\gcd(a,b) = 1$.
\end{proof}

\subsection*{Induction}

\begin{theorem}[Closed form for the first $n$ positive integers]
For every $n \geq 1$,
\[
\sum_{k=1}^{n} k \;=\; \frac{n(n+1)}{2}.
\]
\end{theorem}
\begin{proof}
We argue by induction on $n$, applying Theorem~\ref{thm:induction}.

\emph{Base case.} For $n = 1$ the left-hand side is $\sum_{k=1}^1 k = 1$ and
the right-hand side is $1 \cdot 2 / 2 = 1$. Both equal $1$.

\emph{Inductive step.} Fix $m \geq 1$ and assume the inductive hypothesis
\[
\sum_{k=1}^{m} k = \frac{m(m+1)}{2}. \tag{IH}
\]
We must show the statement for $n = m+1$. Compute
\begin{align*}
\sum_{k=1}^{m+1} k
  &= \left( \sum_{k=1}^{m} k \right) + (m+1) \\
  &\stackrel{(\text{IH})}{=} \frac{m(m+1)}{2} + (m+1) \\
  &= (m+1) \left( \frac{m}{2} + 1 \right) \\
  &= (m+1) \cdot \frac{m+2}{2} \\
  &= \frac{(m+1)\big((m+1)+1\big)}{2},
\end{align*}
which is the claim for $n = m+1$.

By Theorem~\ref{thm:induction}, the identity holds for every $n \geq 1$.
\end{proof}

\section*{Connection to ML}

Every formal guarantee in a machine-learning paper is a quantified proposition:
a convergence rate is $\forall \epsilon > 0,\, \exists N,\, \forall n \geq N,\,
\|\theta_n - \theta^*\| < \epsilon$; a generalisation bound is $\forall \delta,
\Pr[\,\text{risk gap} \leq f(\delta)\,] \geq 1 - \delta$. Identifying the
proof strategy --- direct construction, contrapositive, contradiction, or
induction (often on network depth or training step) --- is the first step in
critiquing the paper. For example, a Lipschitz bound for a deep network is
classically proved by induction on the layer index, with the inductive
hypothesis that the partial network up to layer $k$ has Lipschitz constant
$\prod_{i \leq k} L_i$.

\section*{Exercises}

\begin{enumerate}
  \item (Direct.) Prove that the sum of two odd integers is even.
  \item (Contradiction.) Prove that there is no largest prime. \emph{Hint:} if
        $p_1, \dots, p_n$ were all the primes, consider $N = p_1 \cdots p_n + 1$.
  \item (Induction.) Prove that $\displaystyle \sum_{k=1}^n k^2 = \frac{n(n+1)(2n+1)}{6}$
        for every $n \geq 1$.
  \item (Truth tables.) Verify by truth table that
        $(p \Rightarrow q) \Leftrightarrow (\neg q \Rightarrow \neg p)$
        is a tautology.
\end{enumerate}

\end{document}
